% IEEE-style abstract (180--230 words)
% Rules enforced:
%   - Decision-support framing only; no clinical validation claims.
%   - No FDA approval, no autonomous diagnosis, no radiologist replacement.
%   - All metrics match KNOWN_METRICS.md exactly.

\begin{abstract}
Manual slice-by-slice delineation of the kidney and renal tumors in contrast-enhanced abdominal CT is labor-intensive and subject to substantial inter-observer variability, limiting throughput in clinical workflows.
We present NephroVision, an academic decision-support prototype for automated three-dimensional kidney and renal tumor/cyst segmentation.
The system ingests NIfTI CT volumes, applies Hounsfield-unit clipping ($[-200, 300]$) and min--max normalization, and feeds preprocessed data into a 3D U-Net trained on the KiTS23 challenge dataset.
Inference uses a sliding-window strategy ($64 \times 192 \times 192$ voxels, $50\%$ overlap) with eight-fold test-time augmentation and conservative connected-component post-processing to suppress spurious detections.
Segmentation masks, volumetric statistics, and an interactive browser-based 3D visualization are produced for physician review.
Comprehensive development documentation records preprocessing decisions, training and inference configurations, engineering choices, and tumor/cyst segmentation challenges.
Evaluated on 64 independent held-out KiTS23 test cases, NephroVision achieved a kidney Dice of $0.9307 \pm 0.064$, tumor Dice of $0.6558 \pm 0.262$, mean Dice of $0.7933$, and a tumor detection rate of $100\%$ ($64/64$).
Robust kidney segmentation and reliable tumor detection are demonstrated; tumor boundary precision is moderate, consistent with the inherent difficulty of segmenting small, low-contrast lesions.
The system is not clinically approved and does not replace radiologist interpretation.
\end{abstract}

\begin{IEEEkeywords}
Deep learning, medical image segmentation, kidney tumor detection, computed tomography, decision-support system, 3D U-Net.
\end{IEEEkeywords}
